\chapter{A VERY BRIEF DIGRESSION ON THE LANGUAGE OF CATEGORIES}\label{categories}

In mathematics we study things (objects) and certain mappings between them (morphisms).
To mention just a few, sets and functions, groups and homomorphisms, topological spaces
and continuous maps, vector spaces and linear transformations, and Hilbert spaces and
bounded linear maps.  These examples come from different branches of mathematics---set
theory, group theory, topology, linear algebra, and functional analysis, respectively.
But these different areas have many things in common: in many field terms like product,
coproduct, subobject, quotient, pullback, isomorphism, and projective limit  appear.
Category theory is an attempt to unify and formalize some of these common concepts.
In a sense, category theory is the study of what different branches of mathematics
have in common.  Perhaps a better description is: categorical language tells us ``how things work'',
not ``what they are''.

In these notes, indeed any text at this level, ubiquitously uses the language of sets without
assuming a detailed prior study of axiomatic set theory. Similarly, we will cheerfully use the
\emph{language} of categories without first embarking on a foundationally satisfactory study of
category theory (which itself is a large and important area of research).  Sometimes textbooks make
learning even the language of categories challenging by leaning heavily on one's algebraic background.
Just as most people are comfortable using the language of sets (whether or not they have made a serious
study of set \emph{theory}), nearly everyone should find the use of categorical language both convenient
and enlightening without elaborate prerequisites.

For those who wish to delve more deeply into the subject I can recommend a very gentle entr\'ee to the world
of categories which appears as Chapter 3 of Semadeni's beautiful book~\cite{Semadeni:1971}.  A classic
text written by one of the founders of the subject is~\cite{MacLane:1971}.  A more recent text is~\cite{LawvereS:1997}.

By pointing to unifying principles the language of categories often provides striking insight into
``the way things work'' in mathematics. Equally importantly, one gains in efficiency by not having to go
through essentially the same arguments over and over again in just slightly different contexts.

For the moment we do little more than define ``object'', ``morphism'', and ``functor'', and give a few examples.









\section{Objects and Morphisms}\label{C0151}
\begin{defn} Let $\sfml A$ be a class, whose members we call
 \index{object (of a category)}%
\df{objects}. For every pair $(S,T)$ of objects we associate a class $\mor ST$, whose
members we call
 \index{morphism!of a category}%
\df{morphisms} (or
 \index{arrow}%
\df{arrows}) from $S$ to $T$. We assume that $\mor ST$ and $\mor UV$ are disjoint
unless $S = U$ and $T = V$.

We suppose further that there is an operation $\circ$ (called \df{composition}) that
associates with every $\alpha \in \mor ST$ and every $\beta \in \mor TU$ a morphism
$\beta \circ \alpha \in \mor SU$ in such a way that:
 \begin{enumerate}
  \item $\gamma \circ (\beta \circ \alpha) = (\gamma \circ
\beta) \circ \alpha$ whenever $\alpha \in \mor ST$, $\beta \in \mor TU$, and $\gamma \in
\mor UV$;
  \item for every object $S$ there is a morphism $I_S \in
\mor SS$ satisfying $\alpha \circ I_S = \alpha$ whenever $\alpha \in \mor ST$ and $I_S
\circ \beta = \beta$ whenever $\beta \in \mor RS$.
 \end{enumerate}
Under these circumstances the class $\sfml A$, together with the associated families of
morphisms, is a
 \index{category}%
\df{category}.

We will reserve the notation
 \index{<arrowto@$S \to^\alpha T$ (morphism in a category)}%
$S \to^\alpha T$ for a situation in which $S$ and $T$ are objects in some category and
$\alpha$ is a morphism belonging to $\mor ST$. As is the case with groups and vector
spaces we usually omit the composition symbol $\circ$ and write
  \index{<binopconcat@$\beta\alpha$ (notation for composition of morphisms)}%
  \index{conventions!notation for composition!of morphisms}%
$\beta\alpha$ for $\beta \circ \alpha$.

A category is
 \index{small}%
 \index{category!small}%
\df{locally small} if $\mor ST$ is a set for every pair of objects $S$ and $T$ in the category.  It is
 \index{locally!small}%
 \index{category!locally small}%
\df{small} if, in addition, the class of all its objects is a set.  In these notes the categories in which we are interested
will be locally small.
\end{defn}

\begin{exam} The category
 \index{category!$\cat{SET}$ as a}%
 \index{SET@$\cat{SET}$!the category}%
$\cat{SET}$ has sets for objects and functions (maps) as morphisms.
\end{exam}

\begin{exam} The category
 \index{category!$\cat{AbGp}$ as a}%
 \index{AbGp@$\cat{AbGp}$!the category}%
$\cat{AbGp}$ has Abelian groups for objects and group homomorphisms as morphisms.
\end{exam}

\begin{exam} The category
 \index{category!$\cat{VEC}$ as a}%
 \index{VEC@$\cat{VEC}$!the category}%
$\cat{VEC}$ has vector spaces for objects and linear transformations as morphisms.
\end{exam}







\begin{defn}\label{C008111} Let $\le$ be a relation on a nonempty set $S$.
 \begin{enumerate}
  \item If the relation
 \index{<le@$\le$ (notation for a partial ordering)}%
 $\le$ is reflexive (that is, $x \le x$ for all $x \in S$) and transitive (that is, if $x \le y$ and $y \le z$, then $x  \le z$), it is a
 \index{reflexive}%
 \index{transitive}%
 \index{preordering}%
 \index{ordering!pre-}%
\df{preordering}.
  \item If $\le$ is a preordering and is also antisymmetric (that is, if $x \le y$ and $y \le x$, then $x = y$), it is a
 \index{antisymmetric}%
 \index{partial!ordering}%
 \index{ordering!partial}%
\df{partial ordering}.
  \item Elements $x$ and $y$ in a set on which a preordering has been defined are
 \index{comparable}%
\df{comparable} if either $x \le y$ or $y \le x$.
  \item If $\le$ is a partial ordering with respect to which any two elements are comparable, it is a
 \index{linear!ordering}%
 \index{ordering!linear}%
\df{linear ordering} (or a
 \index{total!ordering}%
 \index{ordering!total}%
\df{total ordering}).
  \item If the relation $\le$ is a preordering (respectively, partial ordering,
linear ordering) on $S$, then the pair $(S,\le)$ is a
 \index{preordered set}%
\df{preordered set} (respectively,
 \index{partially ordered set}%
\df{partially ordered set},
 \index{linearly ordered set}%
\df{linearly ordered set}).
  \item A linearly ordered subset of a partially ordered set $(S,\le)$ is a
 \index{chain}%
\df{chain} in~$S$.
 \end{enumerate}
We may write $b \ge a$ as a substitute for $a \le b$.  The notation
 \index{<l@$<$ (notation for a strict ordering)}%
$a < b$ ( or equivalently, $b > a$) means that $a \le b$ and $a \ne b$.
\end{defn}

\begin{defn} A map $f\colon S \sto T$ between preordered sets is
 \index{order!preserving map}%
\df{order preserving} if $x \le y$ in $S$ implies $f(x) \le f(y)$ in~$T$.
\end{defn}

\begin{exam} The category
 \index{category!$\cat{POSET}$ as a}%
 \index{POSET@$\cat{POSET}$!the category}%
$\cat{POSET}$ has partially ordered sets for objects and order preserving maps as morphisms.
\end{exam}

\begin{notn} Let $f \colon S \sto T$ be a function between sets. Then for $A \subseteq S$ we define
$f^{\sto}(A) = \{f(x)\colon x \in A\}$ and for $B \subseteq T$ we define
 \index{<superscript@$f^{\sto}(A)$ or $f(A)$ (image of $A$ under $f$)}%
 \index{<superscript@$f^{\gets}(B)$ (preimage of $B$ under $f$)}%
$f^{\gets}(B) = \{ x \in S\colon f(x) \in B\}$. We say that $f^{\sto}(A)$ is the
 \index{image}%
\df{image of $A$ under $f$} and that $f^{\gets}(B)$ is the
 \index{preimage}%
\df{preimage} (or the
 \index{inverse!image}%
 \index{image!inverse}%
\df{inverse image}) \df{of $B$ under~$f$}.  Sometimes, to avoid clutter, I will succumb to the dubious practice of
writing $f(A)$ for~$f^{\sto}(A)$.
\end{notn}

\begin{defn} Recall that a family $\sfml T$ of subsets of a set $X$ is a
 \index{topology}%
\df{topology} on $X$ if it contains both the empty set $\emptyset$ and $X$ itself and is closed under taking of unions
($\bigcup\sfml S \in \sfml T$ whenever $\sfml S \subseteq \sfml T$) and of finite intersections ($\bigcap\sfml F \in \sfml T$
whenever $\sfml F \subseteq \sfml T$ and $\sfml F$ is finite).  A nonempty set $X$ together with a topology on $X$ is a
 \index{topological!space}%
 \index{space!topological}%
\df{topological space}. If $(X,\sfml T)$ is a topological space, a member of $\sfml T$ is an
 \index{open}%
 \index{<binrel@$\open UX$ ($U$ is an open subset of~$X$)}%
\df{open set}.  To indicate that a set $U$ is an open subset of $X$ we often write $\open UX$.  The complement of an open set is a
 \index{closed}%
\df{closed set}.  A function $f\colon X \sto Y$ between topological spaces is
 \index{continuous!function}%
 \index{function!continuous}%
\df{continuous} if the inverse image of every open set is open; that is, if $\open{f^{\gets}(V)}X$ whenever $\open VY$.
(If you have not previously encountered topological spaces, look at the first two items in both the first section of
chapter 10 and the first section of chapter 11 of my online notes~\cite{Erdman:2007}.)
\end{defn}

\begin{exam} Every nonempty set $X$ admits a smallest (coarsest) topology; it is the one consisting of exactly two sets,
$\emptyset$ and~$X$.  This is the
 \index{indiscrete}%
 \index{topology!indiscrete}%
\df{indiscrete} topology on~$X$. (Note the spelling: there is another word in the English language, ``indiscreet'',
which means something quite different.)  The set $X$ also admits a largest (finest) topology, which comprises the
family $\sfml P(X)$ of all subsets of~$X$.  This is the
 \index{discrete}%
 \index{topology!discrete}%
\df{discrete} topology on~$X$.
  \end{exam}

\begin{exam} The category
 \index{category!$\cat{TOP}$ as a}%
 \index{TOP@$\cat{TOP}$!the category}%
$\cat{TOP}$ has topological spaces for objects and continuous functions as morphisms.
\end{exam}

\begin{defn}\label{defn_alg} Let $(A,+,M)$ be a vector space over $\K$ which is equipped with
another binary operation $A \times A \sto A$ where $(x,y) \mapsto x \cdot y$ in such
a way that $(A,+,\cdot\,)$ is a ring. (The notation $x \cdot y$ is usually shortened to $xy$.) If additionally the equations
  \begin{equation}\label{eq_def_algebra}
     \alpha (xy) = (\alpha x)y = x(\alpha y)
  \end{equation}
hold for all $x$, $y \in A$ and $\alpha \in \K$, then $(A,+,M,\,\cdot\,)$ is an
 \index{algebra}%
\df{algebra} over the field~$\K$ (sometimes referred to as a \df{linear associative
algebra}).  We abuse terminology in the usual way by writing such things as, ``Let $A$ be an
algebra.'' We say that an algebra $A$ is
 \index{unital!algebra}%
 \index{algebra!unital}%
\df{unital} if its underlying ring $(A,+,\cdot\,)$ has a multiplicative identity; that is, if there exists
 \index{multiplicative!identity}%
 \index{identity!multiplicative}%
an element $\vc 1_A \ne \vc 0$ in $A$ such that $\vc 1_A \cdot x = x \cdot \vc 1_A = x$ for all $x \in A$.
  And it is
 \index{commutative!algebra}%
 \index{algebra!commutative}%
\df{commutative} if its ring is; that is, if $xy = yx$ for all $x$, $y \in A$.

A subset $B$ of an algebra $A$ is a
 \index{subalgebra}%
\df{subalgebra} of $A$ if it is an algebra under the operations it inherits from~$A$.  A
subalgebra $B$ of a unital algebra $A$ is a
 \index{unital!subalgebra}%
 \index{subalgebra!unital}%
\df{unital subalgebra} if it contains the multiplicative identity of~$A$.

\textbf{CAUTION.} To be a unital subalgebra it is \emph{not} enough for $B$ to have a
multiplicative identity of its own; it must contain the identity of~$A$.  Thus, an
algebra can be both unital and a subalgebra of $A$ without being a unital subalgebra
of~$A$.

A map $f\colon A \sto B$ between algebras is an
 \index{homomorphism!of algebras}%
 \index{algebra!homomorphism}%
\df{(algebra) homomorphism} if it is a linear map between $A$ and $B$ as vector spaces which preserves multiplication
(that is, $f(xy) = f(x)f(y)$ for all $x$, $y \in A$).  In other words, an algebra homomorphism is a linear ring homomorphism.
It is a
 \index{unital!algebra homomorphism}%
 \index{algebra!homomorphism!unital}%
 \index{homomorphism!of algebras!unital}%
\df{unital (algebra) homomorphism} if it preserves identities; that is, if both $A$ and $B$ are unital algebras and $f(\vc 1_A) = 1_B$.
The
 \index{kernel!of an algebra homomorphism}%
\df{kernel} of an algebra homomorphism $f\colon A \sto B$ is, of course, $\{a \in A \colon f(a) = \vc 0\}$.

If $f^{-1}$ exists and is also an algebra homomorphism, then $f$ is an
 \index{isomorphism!of algebras}%
\df{isomorphism} from $A$ to $B$.  If an isomorphism from $A$ to $B$ exists, then $A$ and
$B$ are
 \index{isomorphic}%
\df{isomorphic}.
\end{defn}

\begin{exam} The category
 \index{category!$\cat{ALG}$ as a}%
 \index{ALG@$\cat{ALG}$!the category}%
$\cat{ALG}$ has algebras for objects and algebra homomorphisms for morphisms.
\end{exam}

The preceding examples are examples of \emph{concrete categories}---that is, categories
in which the objects are sets (together, usually, with additional structure) and the
morphism are functions (usually preserving this extra structure).  In these notes the
categories of interest to us are both concrete and locally small.  Here (for those who are curious) is a
more formal definition of \emph{concrete category}.

\begin{defn}\label{C015124} A category $\sfml A$ together with a function $\abs{\hphantom{O}}$ which assigns
to each object $A$ in $\sfml A$ a set $\abs A$ is a
 \index{concrete category}%
 \index{category!concrete}%
\df{concrete} category if the following conditions are satisfied:
 \begin{enumerate}
    \item every morphism $A \to^f B$ is a function from $\abs A$ to~$\abs B$;
    \item each identity morphism $I_A$ in $\sfml A$ is the identity function on~$\abs A$; and
    \item composition of morphisms $A \to^f B$ and $B \to^g C$ agrees with composition
of the functions $f\colon \abs A \sto \abs B$ and $g\colon \abs B \sto \abs C$.
 \end{enumerate}
If $A$ is an object in a concrete category $\sfml A$, then $\abs A$ is the
 \index{underlying set}%
 \index{set!underlying}%
 \index{<unop@$\abs A$ (underlying set of the object~$A$)}%
\df{underlying set} of~$A$.
\end{defn}

Although it is true that the categories of interest in these notes are concrete
categories, it may nevertheless be interesting to see an example of a category that is
\emph{not} concrete.

\begin{exam}\label{C015127} Let $G$ be a monoid (that is, a semigroup with identity).
 \index{monoid}%
Consider a category $\cat C$ having exactly one object, which we call~$\star$. Since there
is only one object there is only one family of morphisms $\mor{\star}{\star}$, which we
take to be~$G$. Composition of morphisms is defined to be the monoid multiplication.
That is, $a \circ b := ab$ for all $a$, $b \in G$. Clearly composition is associative and
the identity element of $G$ is the identity morphism. So $\cat C$ is a category.
\end{exam}

\begin{defn} In any concrete category we will call an injective morphism a
 \index{monomorphism}%
\df{monomorphism} and a surjective morphism an
 \index{epimorphism}%
\df{epimorphism}.
\end{defn}

\begin{cau} The definitions above reflect the original Bourbaki use of the term and are the
ones most commonly adopted by mathematicians outside of category theory itself where
``monomorphism'' means ``left cancellable'' and ``epimorphism'' means ``right
cancellable''.  (Notice that the terms \emph{injective} and \emph{surjective} may not
make sense when applied to morphisms in a category that is not concrete.)

A morphism $B \to^g C$ is
 \index{left!cancellable}%
 \index{cancellable!left}%
\df{left cancellable} if whenever morphisms $A \to^{f_1} B$ and $A \to^{f_2} B$ satisfy
$gf_1 = gf_2$, then $f_1 = f_2$.  Mac Lane suggested calling left cancellable morphisms
 \index{monic}%
\df{monic} morphisms.  The distinction between monic morphisms and monomorphisms turns
out to be slight.  In these notes almost all of the morphisms we encounter are monic if
and only if they are monomorphisms. As an easy exercise prove that any injective morphism
in a (concrete) category is monic.  The converse sometimes fails.

In the same vein Mac Lane suggested calling a
 \index{right!cancellable}%
 \index{cancellable!right}%
\emph{right cancellable} morphism (that is, a morphism $A \to^f B$ such that whenever
morphisms $B \to^{g_1} C$ and $B \to^{g_2} C$ satisfy $g_1f = g_2f$, then $g_1 = g_2$) an
 \index{epic}%
\df{epic} morphism.  Again it is an easy exercise to show that in a (concrete) category
any epimorphism is epic.  The converse, however, fails in some rather common categories.
\end{cau}

\begin{defn} The terminology for inverses of morphisms in categories is essentially the same
as for functions.  Let $A \to^\alpha B$ and $B \to^\beta A$ be
morphisms in a category. If $\beta \circ \alpha = I_A$, then $\beta$ is a
 \index{left!inverse!of a morphism}%
 \index{inverse!left}%
\df{left inverse} of $\alpha$ and, equivalently, $\alpha$ is a
 \index{right!inverse!of a morphism}%
 \index{inverse!right}%
\df{right inverse} of $\beta$. We say that the morphism $\alpha$ is an
 \index{isomorphism!in a category}%
\df{isomorphism} (or is
 \index{invertible!morphism}%
\df{invertible}) if there exists a morphism $B \to^\beta A$ which is both a left and a
right inverse for~$\alpha$.  Such a function is denoted by $\alpha^{-1}$ and is called the
 \index{<unopur@$\alpha^{-1}$ (inverse of a morphism $\alpha$)}%
 \index{inverse!of a morphism}%
\df{inverse} of~$\alpha$.
\end{defn}

\begin{prop} Let $A \to^\alpha B$ be a morphism in a category.  If $\alpha$ has both a left inverse $\lambda$ and a
right inverse $\rho$, then $\lambda = \rho$ and consequently $\alpha$ is an isomorphism.
\end{prop}

In any concrete category one can inquire whether every bijective morphism (that is, every
map which is both a monomorphism and an epimorphism) is an isomorphism.  The answer is
often a trivial \emph{yes} (as in $\cat{SET}$, $\cat{AbGp}$, and $\cat{VEC}$) or a trivial
\emph{no} (for example, in the category $\cat{POSET}$
 \index{bijective morphisms!need not be invertible!in $\cat{POSET}$}%
 \index{POSET@$\cat{POSET}$!bijective morphism in}%
 \index{POSET@$\cat{POSET}$!the category}%
of partially ordered sets and order preserving maps.  But on
occasion the answer turns out to be a fascinating and deep result (see for example the \emph{open mapping
theorem}~\ref{C069414}).

\begin{exam} In the category $\cat{SET}$
 \index{bijective morphisms!are invertible!in $\cat{SET}$}%
 \index{SET@$\cat{SET}$!bijective morphism in}%
every bijective morphism is an isomorphism.
\end{exam}

\begin{exam}\label{C015216} The category
 \index{LAT@$\cat{LAT}$!the category}%
 \index{category!$\cat{LAT}$ as a}%
$\cat{LAT}$ has lattices as objects and lattice
homomorphisms as morphisms. (A
 \index{lattice}%
\df{lattice} is a partially ordered set in which every pair of elements has an infimum and a supremum and a
 \index{lattice!homomorphism}%
 \index{homomorphism!lattice}%
\df{lattice homomorphism} is a map between lattices which preserves infima and suprema.)
In this category bijective morphisms
 \index{bijective morphisms!are invertible!in $\cat{LAT}$}%
 \index{LAT@$\cat{LAT}$!bijective morphism in}%
are isomorphisms.
\end{exam}

\begin{exam} A
 \index{homeomorphism}%
\df{homeomorphism} is a continuous bijection $f$ between topological spaces whose inverse $f^{-1}$ is also continuous.  Homeomorphisms
are the isomorphisms in the category $\cat{TOP}$.  In this category not every continuous bijection is a homeomorphism.
\end{exam}

\begin{exam} If in the category $\cat C$ of example~\ref{C015127} the monoid $G$ is a group
then every morphism in $\cat C$ is an isomorphism.
\end{exam}






\section{Functors}
\begin{defn} If $\cat{A}$ and $\cat{B}$ are categories a
 \index{covariant functor}%
 \index{functor!covariant}%
\df{covariant functor} $F$ from $\cat A$ to $\cat B$ (written ${\cat A}\to^F{\cat B}$) is
a pair of maps: an
 \index{object!map}%
 \index{map!object}%
\df{object map} $F$ which associates with each object $S$ in $\cat A$ an object $F(S)$ in
$\cat B$ and a
 \index{morphism!map}%
 \index{map!morphism}%
\df{morphism map} (also denoted by $F$) which associates with each morphism $f \in \mor
ST$ in $\cat A$ a morphism $F(f) \in \mor {F(S)}{F(T)}$ in $\cat B$, in such a way that
 \begin{enumerate}
  \item $F(g \circ f) = F(g) \circ F(f)$ whenever $g \circ f$ is defined in $\cat A$; and
  \item $F(\id S) = \id{F(S)}$ for every object $S$ in $\cat A$.
 \end{enumerate}

The definition of a
 \index{contravariant functor}%
 \index{functor!contravariant}%
\df{contravariant functor} ${\cat A}\to^F{\cat B}$ differs from the preceding definition
only in that, first, the morphism map associates with each morphism $f \in \mor ST$ in
$\cat A$ a morphism $F(f) \in \mor {F(T)}{F(S)}$ in $\cat B$ and, second, condition (a)
above is replaced by
 \begin{enumerate}
  \item[(\rm{a'})] $F(g \circ f) = F(f) \circ F(g)$ whenever $g \circ f$
is defined in~$\cat A$.
 \end{enumerate}
\end{defn}

\begin{exam}\label{functors005exam} A
 \index{forgetful functor}%
 \index{functor!forgetful}%
\df{forgetful functor} is a functor that maps objects and morphisms from a category $\cat
C$ to a category $\cat C'$ with less structure or fewer properties.  For example, if $V$
is a vector space, the functor $F$ which ``forgets'' about the operation of scalar
multiplication on vector spaces would map $V$ into the category of Abelian groups. (The
Abelian group $F(V)$ would have the same set of elements as the vector space $V$ and the
same operation of addition, but it would have no scalar multiplication.)  A linear map
$T\colon V \sto W$ between vector spaces would be taken by the functor $F$ to a group
homomorphism $F(T)$ between the Abelian groups $F(V)$ and~$F(W)$.

Forgetful functor can ``forget'' about properties as well.  If $G$ is an object in the
category of Abelian groups, the functor which ``forgets'' about commutativity in Abelian
groups would take $G$ into the category of groups.

It was mentioned in the preceding section that all the categories that are of interest in
these notes are concrete categories (ones in which the objects are sets with additional
structure and the morphisms are maps which preserve, in some sense, this additional
structure). We will have several occasions to use a special type of forgetful
functor---one which forgets about all the structure of the objects except the underlying
set and which forgets any structure preserving properties of the morphisms. If $A$ is an
object in some concrete category $\cat C$, we denote by
 \index{<unop@$\abs A$ (the forgetful functor acting on~$A$)}%
$\abs{A}$ its underlying set. And if \mbox{$A~\to^f~B$} is a morphism in $\cat C$ we
denote by $\abs f$ the map from $\abs A$ to $\abs B$ regarded simply as a function
between sets.  It is easy to see that $\abs{\hphantom{M}}$\,, which takes objects in
$\cat C$ to objects in $\cat{SET}$ (the category of sets and maps) and morphisms in $\cat
C$ to morphisms in $\cat{SET}$, is a covariant functor.

In the category $\cat{VEC}$ of vector spaces and linear maps, for example,
$\abs{\hphantom{M}}$ causes a vector space $V$ to ``forget'' about both its addition and
scalar multiplication ($\abs V$ is just a set). And if $T\colon V \sto W$ is a linear
transformation, then $\abs T\colon \abs V \sto \abs W$ is just a map between sets---it
has ``forgotten'' about preserving the operations.
\end{exam}


\begin{defn} A partially ordered set is
 \index{order!complete}%
 \index{complete!order}%
\df{order complete} if every nonempty subset has a supremum (that is, a least upper
bound) and an infimum (a greatest lower bound).
\end{defn}

\begin{defn} Let $S$ be a set.  Then the
 \index{power set}%
 \index{P@$\sfml P(S)$ (power set of $S$)}%
\df{power set} of $S$, denoted by $\sfml P(S)$, is the family of all subsets of~$S$.
\end{defn}

\begin{exam}[The power set functors] Let $S$ be a nonempty set.
 \begin{enumerate}
  \item The power set $\sfml P(S)$ of $S$ partially ordered by $\subseteq$ is
order complete.
  \item The class of order complete partially ordered sets and order
preserving maps is a category.
 \index{functor!power set}%
 \index{power set!functor}%
  \item For each function $f$ between sets let $\sfml P(f) = f^{\sto}$. Then
$\sfml P$ is a covariant functor from the category of sets and functions to the category
of order complete partially ordered sets and order preserving maps.
  \item For each function $f$ between sets let $\sfml P(f) = f^{\gets}$. Then
$\sfml P$ is a contravariant functor from the category of sets and functions to the
category of order complete partially ordered sets and order preserving maps.
 \end{enumerate}
\end{exam}

\begin{defn}\label{defn_vsp_adjoint} Let $T\colon V \sto W$ be a linear map between vector spaces. For every
$g \in W^{\#}$ let $T^{\#}(g) = g\,T$.  Notice that $T^{\#}(g) \in V^{\#}$.  The map $T^{\#}$ from the
vector space $W^{\#}$ into the vector space~$V^{\#}$ is the (vector space)
 \index{adjoint!vector space}%
 \index{vector!space!adjoint map}%
 \index{<unopur@$T^\#$ (vector space adjoint)}%
\df{adjoint} map of~$T$.
\end{defn}

\begin{exam} The pair of maps $V \mapsto V^{\#}$ (taking a vector space to its algebraic dual) and
$T \mapsto T^{\#}$ (taking a linear map to its dual) is a contravariant functor from the category $\cat{VEC}$
of vector spaces and linear maps to itself.
\end{exam}

\begin{exam}\label{0007606} If $\cat C$ is a category let
 \index{category@$\cat{C^2}$ (pairs of objects and morphisms)}%
$\cat{C^2}$ be the category whose objects are ordered pairs of objects in $\cat C$ and
whose morphisms are ordered pairs of morphisms in~$\cat C$.  Thus if $A \to^f C$ and $B
\to^g D$ are morphisms in $\cat C$, then $(A,B) \to^{(f,g)} (C,D)$ (where $(f,g)(a,b) =
\bigl(f(a),g(b)\bigr)$ for all $a \in A$ and $b \in B$) is a morphism in~$\cat{C^2}$.
Composition of morphism is defined in the obvious way: $(f,g) \circ (h,j) = (f \circ h, g
\circ j)$. We define the
 \index{diagonal!functor}%
 \index{functor!diagonal}%
\df{diagonal functor} $\cat C \to^{\ftr D} \cat{C^2}$ by $\ftr D(A) := (A,A)$. This is a
covariant functor.
\end{exam}

\begin{exam}\label{C029717} Let $X$ and $Y$ be topological spaces and $\phi\colon X \sto Y$
be continuous. Define $\fml C \phi$ on $\fml C(Y)$ by
    \[ \fml C \phi(g) = g \circ \phi \]
for all $g \in \fml C(Y)$.  Then the pair of maps $X \mapsto \fml C(X)$ and
 \index{C@$\fml C$ (the functor)}%
$\phi \mapsto \fml C \phi$ is a contravariant functor from the category $\cat{TOP}$ of topological
spaces and continuous maps to the category of unital algebras and unital algebra
homomorphisms.
\end{exam}













\endinput

